\documentclass[10pt,a4paper]{amsart}
\usepackage[margin=19mm]{geometry}
\usepackage{amsmath,amssymb,mathtools}
\usepackage[scr=rsfs,cal=euler]{mathalfa}
\usepackage{xfrac}
\usepackage[unicode,hidelinks,bookmarks=false]{hyperref}
\newcommand{\ie}{i.e.\ }
\newcommand{\cf}{cf.\ }
\newcommand{\resp}{respectively\ }
\newcommand{\Subsection}{\subsection}
\providecommand{\nobreakdash}{\nobreak\hbox{-}}
\setlength{\emergencystretch}{2em}
\multlinegap0pt
\usepackage{amssymb}
\newtheorem{prop}{Proposition}
\title{Rapid Decay Subalgebras of C$^*$-Algebras}
\author{Shelley Hebert, Slawomir Klimek, Matt McBride}
\date{Source: arXiv:2605.08977v1 (CC BY 4.0)}
\begin{document}
\maketitle

\noindent\emph{Source and licence.} Rapid Decay Subalgebras of C$^*$-Algebras. Version arXiv:2605.08977v1 (CC BY 4.0), distributed by arXiv under the Creative Commons Attribution 4.0 International licence, http://creativecommons.org/licenses/by/4.0/, as recorded in the arXiv licence metadata for this exact version and shown on the abstract page https://arxiv.org/abs/2605.08977v1. Full licence notice: \texttt{licenses/arxiv-2605.08977-CC-BY-4.0.txt}.\par\smallskip

\noindent\textit{Editorial scope.} Counted unit: the article's prop pnxyprop (source0.tex lines 1046--1054). The article's complete original proof of it is delivered with it (source0.tex lines 1056--1097, one \texttt{proof} environment of 42 lines). No mathematics is written, repaired or re-derived: the statement and the proof are the article's own lines, in the article's own notation, and no retained line is substituted or re-flowed.  No further source-local context is needed: every cross-reference of every retained line resolves inside the two delivered spans.  Nothing was omitted from any retained span, so the package carries no omission record.  No retained line contains a citation, so the wrapper carries no bibliography and no bibliography entry.  No retained line contains a figure environment, an image-inclusion command or drawing code, so no image is delivered and the package declares no asset dependency.  Editorial formatting only, in addition: an AMS \texttt{amsart} preamble that restates the article's own declarations and macros used by the retained lines, namely the environments \texttt{prop} on separate counters, together with the standard packages those lines need.  The retained environments print in this document as Proposition 1 in order of appearance, whereas the article prints the same items with the numbering of its own full document.  The complete native source of the article as distributed in the arXiv e-print for this exact version is bundled unchanged as \texttt{sources/200189/source0.tex}.
\begin{prop}\label{pnxyprop}
With the above notation, the operators $P_{(n,x,y)}$ have the following properties for all  $n=0,1,2, \ldots$ and $0\leq x,y,z,w<s_n$:
\begin{enumerate}
\item $P_{(n,x,y)}^*=P_{(n,y,x)}$
\item  $P_{(n,x,y)}P_{(n,w,z)}=\delta_{y,w}P_{(n,x,z)}$
\item  $\sum\limits_{x=0}^{s_n-1}P_{(n,x,x)}=I$
\item   $\sum\limits_{j=0}^{s_{n+1}/s_n-1}P_{(n+1,x+js_n,y+js_n)}=P_{(n,x,y)}$
\end{enumerate}
\end{prop}

\begin{proof}
The first item follows the observation that the inner product $(E_l,P_{(n,x,y)}E_k)$ is not zero only if $l=k+x-y$ and 
$s_n|(k-y)$, which is equivalent to $k=l+y-x$ and $s_n|(l-x)$.

For the second item, consider the following calculation:
\begin{equation*}
\begin{aligned}
P_{(n,x,y)}P_{(n,w,z)}E_k &= P_{(n,x,y)}\left\{
\begin{aligned}
& E_{k+w-z}  &\textrm{ if }s_n|(k-z)\\
&0 &\textrm{else} 
\end{aligned}\right. \\
&= \left\{
\begin{aligned}
& E_{k+w-z+x-y}  &\textrm{ if }s_n|(k-z)\textrm{ and }s_n|(k+w-z-y)\\
&0 &\textrm{else.} 
\end{aligned}\right.
\end{aligned}
\end{equation*}
Since $s_n|(k-z)$, it follows that $s_n|(w-y)$, but since $0\le w,y<s_n$, it must be that $w=y$ to have a nonzero outcome.  In other words
\begin{equation*}
P_{(n,x,y)}P_{(n,w,z)} = \delta_{wy}P_{(n,x,z)}\,.
\end{equation*}

For the third item notice that $P_{(n,x,x)}E_k = E_k$ only if $s_n|(k-x)$, otherwise it is zero.  Thus
\begin{equation*}
\sum_{x=0}^{s_n-1}P_{(n,x,x)}E_k = E_k = IE_k\,,
\end{equation*}
because there is only one nonzero term in the sum, namely when $x = k\textrm{ mod }s_n$.

For the last item, notice that
\begin{equation*}P_{(n+1,x+js_n,y+js_n)}E_k= E_{k+x-y}
\end{equation*}
only if $s_{n+1}|(k-y-s_n)$ and zero otherwise.   Like in the previous item, there is only one term in the sum below that produces a nonzero term, namely when $j=(k-y)/s_n\textrm{ mod }s_{n+1}/s_n$.  It follows that
\begin{equation*}
\sum_{j=0}^{s_{n+1}/s_n-1}P_{(n+1,x+js_n,y+js_n)}E_k=\left\{
\begin{aligned}
& E_{k+x-y} &&\textrm{ if }s_n|(k-y) \\
&0 &&\textrm{ else}
\end{aligned}\right. = P_{(n,x,y)}E_k\,.
\end{equation*}
\end{proof}

\end{document}
